\documentclass[letterpaper,12pt]{article}
\usepackage{helvet}
\usepackage{fancyheadings}
\usepackage{verbatim}
\usepackage{hyperref}
\pagestyle{fancy}
\usepackage{graphicx}
\setlength\textwidth{6.5in}
\setlength\textheight{8.5in}
\begin{document}
\title{Literate Programming in the Era of LLMs}
\author{Ramesh Subramonian }
\maketitle
\thispagestyle{fancy}
\lhead{}
\chead{}
\rhead{}
\lfoot{}
\cfoot{}
\rfoot{{\small \thepage}}


\newtheorem{problem_statement}{Problem Statement}
\newtheorem{notation}{Notation}
\newtheorem{convention}{Convention}
\newtheorem{assumption}{Assumption}

\newcommand{\TBC}{\framebox{\textbf{TO BE COMPLETED}}}
\newcommand{\be}{\begin{enumerate}}
\newcommand{\ee}{\end{enumerate}}
\newcommand{\bi}{\begin{itemize}}
\newcommand{\ei}{\end{itemize}}
\newcommand{\bd}{\begin{description}}
\newcommand{\ed}{\end{description}}

\newcommand{\True}{{\tt true} }
\newcommand{\False}{{\tt false} }
\newcommand{\uint}{{\tt uint32\_t} }
\newcommand{\qt}{{\tt qtype\_t} }
\newcommand{\jt}{{\tt jtype\_t} }
\newcommand{\qdf}{{\tt qdf\_rec\_type} }

\newcommand{\jarr}{{\tt j\_array} }
\newcommand{\jobj}{{\tt j\_object} }

\newcommand{\getq}{{\tt get\_qtype}}
\newcommand{\getj}{{\tt get\_jtype}}
\newcommand{\getn}{{\tt get\_length}}
\newcommand{\getrv}{{\tt get\_read\_arr}}
\newcommand{\getwv}{{\tt get\_write\_arr}}
\newcommand{\getnk}{{\tt get\_num\_keys}}
\newcommand{\getrnn}{{\tt get\_read\_nulls}}
\newcommand{\getwnn}{{\tt get\_write\_nulls}}
\newcommand{\chk}{{\tt chk\_qdf}}

\newcommand{\setq}{{\tt set\_qtype}}
\newcommand{\setj}{{\tt set\_jtype}}
\newcommand{\setn}{{\tt set\_length}}

\newcommand{\mkUA}{{\tt make\_uniform\_array}}
\newcommand{\mknnUA}{{\tt make\_uniform\_array\_with\_nulls}}

\newcommand{\setv}{{\tt set\_arr}}
\newcommand{\setnn}{{\tt set\_nn}}

\newcommand{\Q}{{\tt Q}}
\newcommand{\QDF}{{\tt QDF}}


\abstract{LLMs have taken the programming world by storm. 
How does this change the role of the software engineer? 
Assuming that the cost of writing code will approach zero,
what does a software engineer do? 
We posit that there will be some fundamental changes in the way we will 
build software. (1) We will focus more on writing 
unambiguous, semi-formal specifications than we have traditionally done
(2) We will deploy more custom code, tailored to the specific need of the
application, thereby reducing code bloat
(3) We will adopt just-in-time manufacturing approach to software deployment 
In this paper, we use an example to demonstrate how these principles 
are put into practice. 

}

\section{Introduction}
\label{intro}

Traditionally, the cost of writing software 
(and modifying it when new needs arose) was so high that the design of 
code optimized for maintainability. Most seasoned software engineers will have
experienced the sinking feeling when asked to add new requirements to a large
code base, the fear of not fully understanding whether a fix here 
would cause a blow out there. We would argue that assuming 
zero cost of (bug-free) software development flips this paradigm on 
its head. This moves the burden from refactoring code to making sure that the specifications are
unambiguous and consistent and meet the business requirement precisely.

Good software engineering was never about churning out lines of code. 
To quote Lattner, {\em Writing code has never been the goal. Building meaningful
software is.} \cite{CCC_Lattner_2026}.
Instead, it was about designing systems that met the requirements.
Djkstra's claim was never more relevant than it is today --- 
``{\em Computer science is no more about computers than astronomy is about
telescopes.}'' Writing efficient code was a means to an end.
There have been a series of improvements in that process, 
LLMs representing the
latest and admittedly most unexpected quantum jump in efficiency.
Writing formal specifications is a declaration of intent. 
Until the singularity arrives, specifications represent the safest 
bet that we will get what we want when we no longer control the 
translation from intent to a higher level programming language,
en route to a further translation to machine instructions.


Which brings us to the central question that this paper seeks to answer.
\begin{center}
{\em How much code should a coder code if Claude Code could code code?}
\end{center}

For those with little patience to slog through the rest, here's the gist
\be
\item control the data, the format and inter-relationships
\item write specifications for new operators and have the LLM write code in a
  high level programming language which you can review 
\item deploy functionality as needed, at the cost of a small 
  delay when you ask for something new
  \ee

Throwing down the gauntlet to the godfather of ``vibe coding'', 
I wonder why Karpathy chose to write C code (excellent code though it is)
for Llama 2 instead of just writing a set of math formulas
\footnote{\url{https://github.com/karpathy/llama2.c}}.

\subsection{Organization}
In Section~\ref{motivation}, we present the example used to 
ground the claims made above. 
In Section~\ref{designing_for_llms}, we describe how 
our approach to software development changed fundamentally with the inclusion of
LLMs.
In Section~\ref{concrete}, we make the approach and philosophy concrete by 
delving into the details of the implementation.

\section{Motivation}
\label{motivation}

Earlier, 
the author had built a column-oriented analytical database called \Q. 
As with most such databases, greater efficiencies are possible
when data is created and accessed within the database
engine itself. 
However, a database that doesn't share its data with other applications isn't of
much use. This necessitated the development of \QDF\ (\Q's data format) for data
exchange, not for bulk import/export.

We considered several other formats, starting with the usual suspects,
 CSV and JSON. For efficiency reasons, we narrowed our focus to binary formats, like BSON
 and AVRO. Ultimately, we decided it best to roll our own format based on the
 following requirements and usage patterns
 \be
 \item ease of adding fixed-width custom types
 \item optimizing computations on the data without moving or
   transforming it
 \item allowing for in-situ operations,
   although most operations are write-once.
 \item optimizing for the case where data is read more 
   often than it is written. 
 \item reluctance to depend on software which was unashamed to admit that ``{\em The documentation is
   rather sparse right now but we’ll be adding more information
   soon.}''\footnote{\url{https://avro.apache.org/docs/1.12.0/api/c/}}
   \ee

The development work could be broadly divided into two parts
\be
\item Definition of the data format and writing helper functions which allowed
  the creation and modification of data packets in \QDF\ format and access to
  elements of such a packet
\item Definition and development of operators that read from (and occasionally
  wrote to) the data packet.
  A simple example of an operator is a dot product which
accesses two numeric arrays of equal length and returns a number. 
  \ee
Over a period of four years, we found that the changes in
the underlying data format were both infrequent and minor.
As a consequence, the work on getters and setters was minimal. 
Instead, most of the development work involved
\be
\item writing new operators. 
\item optimizing the performance of existing operators 
  \ee
This led to the following questions
\be
\item Could the use of LLMs have reduced the development and maintenance cost?
\item If so, what would the software architecture of such a system look like?
  \ee

Giving GenAI the benefit of the doubt, we assumed the answer to the first
question would be Yes and proceeded to answer the second question. 

\section{Designing for LLMs}
\label{designing_for_llms}

\subsection{Choice of Language}

We chose to build our system using C and Lua and \LaTeX.  
There is nothing sacrosanct about these choices. Other 
choices of compiled language and interpreted
language and document formatting language are likely to work just as well. 
The reasons for this choice were
\be
\item The Ousterhout dichotomy. This is a principle of software design that I
  have found very useful. For those unfamiliar with the idea, 
``{\em Ousterhout's dichotomy claims that there are two general categories of programming languages:
\be
\item low-level, statically-typed systems languages
\item high-level, dynamic scripting languages
  \ee
Systems languages are best for efficiently handling large quantities of data,
implementing algorithms, and can handle significant internal complexity.
Scripting languages are best for glueing other programs together, exploring a
problem, and extending applications.}
''
\footnote{\url{https://gist.github.com/nat-418/bdccb2428cde750c9f42bf9bbc1a50d3}}

\item Completeness. 
It was very much our intent to
have a fully functional system of moderate complexity which would allow us to
suss out any gotchas that a ``Hello World'' program might have glided by.
The author is sufficiently conversant in these languages to meet this goal.

\item Ease of exposition. 
  The simple syntax of C and Lua make the code snippets widely
  readable
\item Ease of implementation.
  ``{\em Lua was designed since its inception to interoperate with
other languages, both by extending --- allowing Lua code to call functions
written in a foreign language --- and by embedding --- allowing foreign code to call functions written in Lua. Lua is
thus implemented not as a standalone
program but as a library with a C API.
This library exports functions that cre ate a new Lua state, load code into a
state, call functions loaded into a state,
access global variables in a state, and
perform other basic tasks. The stand alone Lua interpreter is a tiny application
written on top of the library.}'' \cite{Lua2019}
\item \LaTeX, when combined with {\tt make}, allows us to ``compile''
  documentation from disparate sources. It allows us to create 
  different targets from common sources. 
\ee



\subsection{Control the Data}
We took our inspiration from the quote attributed to Linus Torvalds --- 
{\em ``Bad programmers worry
about the code. Good programmers worry about data structures and their
relationships.} Hence, we focused our efforts on designing an efficient data
format  (Section~\ref{data_format_defn}) and 
and defining how it would be accessed (Section~\ref{accessor_defn}).
This information is passed to the LLM in the form of a \LaTeX\ file. 
For human readability, some scripted transformations are performed. For example,
a Lua table of the form 
\begin{verbatim}
return { Q0 = "void", I32 = "int32_t", FP32 = "float", FP64 = "double", }
\end{verbatim}
is formatted as Table~\ref{sample_q2c} using a Lua script invoked by {\tt make}.
\begin{table}[hb]
\centering
\begin{tabular}{|l|l|} \hline \hline 
{\bf Q Type } & {\bf C Type } \\ \hline \hline
 Q0 & \verb+ void + \\ \hline
 I32 & \verb+ int32_t + \\ \hline
 FP32 & \verb+ float + \\ \hline
 FP64 & \verb+ float + \\ \hline
\hline
\end{tabular}
\caption{Qtypes to C Types}
\label{sample_q2c}
\end{table}

\subsubsection{Define the Data Format}
\label{data_format_defn}

We wanted the data format to resemble JSON but to have efficient 
representations of commonly used types e.g., dataframes. Consider the JSON type
{\tt number}. This is stored in QDF as a 16 byte packet where
\be
\item Byte 0 contains JSON type information --- {\tt number} in this example.
\item Byte 1 contains data type information e.g., {\tt int32\_t} or {\tt double}
\item Byte 2 is the successor offset. While this is usually 0, it helps tell us
  where the successor QDF, if any, is located. This is used when we want to
  make sure that certain data structures are memory aligned, helpful for vector
  instructions.
\item Byte 3 is unused.
\item Bytes 4 to 7 contain an offset to the parent packet, if any. 
  For example,
  in the JSON array \verb+[ true, "abc", 123 ]+, the number 123 needs to 
  know the start of the parent array of which it is a part. Similarly, the
  successor offset in Byte 3 tells us how much padding exists between the
  boolean {\tt true} and the string {\tt "abc"}.
\item Bytes 8 to 15 contain the actual value
  \ee
\subsubsection{Define the Accessors}
\label{accessor_defn}

After defining the data format, we implemented a number of helper 
functions which broadly fell into these categories
\be
\item readers or getters e.g., calling {\tt get\_jtype()} on a QDF packet tells
  us what JSON type it is i.e., {\tt nil, number, string, \ldots}
\item writers or setters
\item makers or constructors e.g., calling {\tt make\_number()} with the 
  type and value of the number to be created as arguments, 
\item destructors were largely handled by Lua's garbage collection
\item checkers, which validated the syntactic structure of a data packet.
  These can be thought of as invariants, 
  inspired by \cite{hoare1969axiomatic}.

  \ee

\subsubsection{Test the Accessors}
Once the accessors were in place, it was necessary to write unit tests which
would give us confidence in their implementation. We started by hand writing
these tests but quickly transitioned as described in
Section~\ref{accessor_test}. %% TODO 

\subsubsection{Define the Operators}

We took our inspiration from Knuth: ``{\em Let us change our traditional attitude to the construction of programs: 
Instead of imagining that our main task is to 
instruct a computer what to do, let us concentrate rather on explaining
to human beings what we want a computer to do.}'' \cite{knuth1984literate}.

Taking this last sentence to its logical conclusion, we chose to 
concentrate {\bf entirely} on the human explanation, 
leaving the job of instructing the computer what to do to the LLM.

In this paper, Knuth argues that the system he proposes, 
{\em WEB, is 
itself is chiefly a combination of two other languages: (1) a document formatting language and (2) a
programming language. My prototype WEB system uses
\TeX\ as the document formatting language and PASCAL
as the programming language, but the same prin ciples would apply equally well if other languages were
substituted. Instead of \TeX, one could use a language
like Scribe or Troff; instead of PASCAL, one could use
ADA, ALGOL, LISP, COBOL, FORTRAN, APL, C, etc.,
or even assembly language. The main point is that WEB
is {\bf inherently bilingual}, and that such a combination of
languages proves to be much more powerful than either
single language by itself. WEB does not make the other
languages obsolete; on the contrary, it enhances them.
}

In our case, we chose to use \LaTeX\ for specifications 
and have the LLM write the code in Lua, C and ISPC \cite{ISPC}
The author's familiarity with these languages allowed for code review.
If LLMs live up to the trajectory claimed for them, 
we expect to gain confidence in them and reduce the need for oversight.

\section{Making it concrete}
\label{concrete}

Let us assume that having defined the data format, we wish to add 
an operator {\tt coalesce}
\footnote{\url{https://www.rdocumentation.org/packages/dplyr/versions/1.0.10/topics/coalesce}}.
We walk the reader through the process, step by step, explaining what the
developer does and what the system does in response


\subsection{Starting Out}
\begin{verbatim}
lQDF = require 'lQDF'
\end{verbatim}
This informs the system about the data definitions and data accessors.
To those familiar with LuaJIT \cite{LuaJIT}, we have
\be
\item {\tt ffi.cdef}'d the function declarations
\item {\tt ffi.load}'d the library that implements the accessors
  \ee

  \subsection{Registering an operator}
\begin{verbatim}
local X = require 'coalesce'
lQDF.Q.register('coalesce', X)
\end{verbatim}

In the first line, we define {\tt X}, the specification of the operator, 
Section~\ref{coalesce_spec}. It consists of the following
\be
\item \verb+generic_tex_spec_file+: \LaTeX\ file containing 
  type-agnostic specifications.  
  See Section D (captioned {\bf Coalesce}) of Appendix for an example
\item \verb+get_subs+: a function which produces substitutions needed
  to specialize the above \LaTeX\ file based on the types of the arguments
\item \verb+lua_calling_spec+: table containing specification used to 
  create a function {\tt
  run()} which is used at runtime to create (if necesssary) 
  and execute a C function
  \ee

In the second line, the call to {\tt Q.register} causes the following
\be
\item Records the information provided
\item Dynamically generates the code for a function {\tt run()}
  using \verb+lua_calling_spec+.
  Currently, this is done by a human authored script. 
  Future work could see this be passed to the LLM as well. 
\item Dynamically compiles the code generated above using 
  Lua's {\tt loadstring()} into a variable of type {\tt function}
\item Adds {\tt coalesce} as a key to the table {\tt lQDF}, 
  the value of which is the function {\tt run}
  \ee

\subsubsection{Coalesce specification}
\label{coalesce_spec}
\verbatiminput{doc_coalesce.lua}

  \subsection{Calling an operator}
  \begin{verbatim}
local x = lQDF({ 1, 2, 3, }
local y = lQDF({ 4, 5, 6, })
local z = lQDF.coalesce(x, y)
  \end{verbatim}

The first two lines create {\tt x, y}, which are objects of type lQDF.
The third line creates {\tt z}, also an object of type {\tt lQDF}. 
Let's take a peek under the covers and see what the system does
to create {\tt z}.
\be
\item Based on the types of the input arguments, {\tt x, y}, 
determine whether a specialized function for the job exists. 
For now, let's assume it doesn't.
\item Make substitutions in the generic \LaTeX\ file e.g., this is where we 
  would find that the {\tt qtype} of {\tt x,y} is {\tt FP64}, as 
  a result of which we would need to create a function whose name is 
  \verb+coalesece_FP64_FP64+
\item Assemble a larger \LaTeX\ file which contains 
  context which does not change across operators. 
  This is the information in Sections~\ref{data_format_defn} and 
  ~\ref{accessor_defn}.
\item Use {\tt pdflatex} to create a PDF with detailed instructions for the LLM.
  See the Appendix for an example of what is sent to the LLM. Note that the
  developer writes only the portion of this document specific to their 
  operator
\item POST the PDF to the LLM. We also post relevant {\tt .h} (created by the
  system architect not the operator writer) that 
  define {\tt structs, enums, macros}, such as 
  \be
\item  \verb+qdf_struct.h+ 
\item  \verb+qtype.h+ 
\item  \verb+jtype.h+ 
\item  \verb+macros.h+ 
  \ee
We instruct the LLM to return 
  \be
\item the function declaration {\tt .h} file 
\item the function definition {\tt .c} file 
\item commands needed to create a {\tt .o} file and a {\tt .so} file. 
 The LLM gives us information such as
  \be
\item compiler to be used e..g, {\tt gcc} or {\tt ispc}
\item include paths
\item libraries that may need to be linked
\item compiler/linker flags 
  \ee
  \ee
\item Execute above commands to create a {\tt .so} file. 
\item \verb+ffi.cdef+ the {\tt .h} file and \verb+ffi.load+ the {\tt .so} file
  created above
\item The previous steps are skipped over if they had been executed before. 
The system remembers past work. Coming across a new symbol which necessitates
the above work can be thought of as a cache miss. 
This is an acceptable tradeoff {\bf if}
\be
\item these cache misses are infrequent
\item after an initial warm-up phase, the system settles into a steady state
\item the deployment is not sensitive to the latency they introduce,
  \ee
\item Provide C access to the Lua objects {\tt x, y, z}
\item Execute the C function using the pointers provided above as arguments
\item Return the Lua objects specified in {\tt out\_args} in
  Section~\ref{coalesce_spec}.
  \ee

\subsection{Saving and Restoring}

What happens when a session ends? If we wish to save operators created in this
session for use in subsequent sessions, we execute \verb+lQDF.Q.save()+. This
records all necessary artifacts on the file system e.g., \LaTeX\ files, pdf
files, {\tt .so files, \ldots }

When we start a new session, we do not need to start from a blank slate.
Instead, we can pick up from where we left off as show below.
\begin{verbatim}
lQDF.Q.restore()
local z = lQDF.coalesce(lQDF({ 1, 2, 3, }), lQDF({ 4, 5, 6, }))
\end{verbatim}

\section{Discussion}

I started this journey as a curious sceptic. I reached this milestone 
as a cautious optimist. I was amazed at the quality of the code produced. 
But I didn't leave me with a sense of dread that something I took pride 
in could be replicated so easily. Instead, it felt liberating. 
I had found the L3 engineer that I had always wanted to work with. 
Maybe tomorrow I will find the colleague I always wanted to pair program with. 

\section{Conclusion}

The rapid pace of progress in AI makes prognostication difficult. 
Does one ignore it and wait for the bubble to burst? 
Does one embrace it entirely like a product manager of mine who wondered why it
took me so long to convert his PowerPoint wireframes into working code?
This paper represents a measured adoption of the seemingly miraculous powers of
AI which we suspect might be relevant for the next few years.

We agree with Lattner's assessment that
{\em As writing code is becoming easier, designing software becomes more
important than ever. As custom software becomes cheaper to create, the real
challenge becomes choosing the right problems and managing the resulting
complexity.  } \cite{CCC_Lattner_2026}.

\subsection{Speculation}

\be
\item Has AI, for the software engineering domain, become a commodity? 
  For the simple example cited in this paper, I could not find any
  appreciable difference in code quality across several LLMs --- ChatGPT,
  DeepSeek, Gemini, \ldots

One could argue that that was because the context provided was concise and
precise. But isn't that what we would want anyway? 
Rather than treating code as all important and 
documentation and comments as annoyances, it seems right to focus effort on
communicating intent. 

So, while more recent models are defying belief by solving Math Olympiads, 
does one need that power in the next few years to produce high 
quality software economically? I suspect not. 

\item Benioff's theatrics in 2000 declaring that ``software is dead'' 
  may well have come home to roost. Why would I pay for software when what I am
  really paying for is specifications? I believe the time has
  come for disciplined software engineering that uses LLMs for code 
  generation to disrupt the enterprise software market. Customizing a system
  for different businesses will not require the same effort that it has historically required. 
  

\item Many software consulting houses 
  (here's looking at you Infosys, TCS, \ldots) will suffer but not because 
  the need for software will disappear. 
  Instead, they have hired, trained and promoted for vocational skills. 
  They are unprepared to pivot to a realm where engineering skills dominate.
  Many universities in India are really trade schools masquerading as 
  bastions of higher learning.  I do {\bf not} use the term ``trade school''
  disparagingly.  I think they have a valuable purpose. 
  I do hope that the mechanic who tunes up my car has attended
  one of those. But that doesn't mean that I'd hire her to design my dream car. 

\item Would there soon come a time where a software engineer 
  apologizes for writing code instead of a specification with the quote various
  atttibuted to Cicero, Pascal, Mark Twain
\begin{center}
  {\em I didn't have time to write a short letter, so I wrote a 
  long one instead.}
  \end{center}

\item Would LLMs change the way we design programming languages? 
So far, the target of human effort has been code for a high-level 
programming language. How would a specification language differ from 
what we have traditionally thought of as a programming language?

\item Over the years, I've gotten a lot of flak for my unorthodox choice of
  programming languages --- C and Lua. But if it is the LLM writing the code,
  not me, does it matter? 
  If I have time, I'd like to re-implement this project so that a 
  Python programmer would just have to {\tt import foo} and start using {\tt
  foo.coalesce()} without caring about the underlying details. This is similar
  to the fact that most {\tt numpy} users don't care or know that the library is
  written in C.

\item Many years ago, I graduated in the middle of a painful recession. A friend
  of mine persuaded a hiring manager at Oracle to take a call from me. I picked
  up the handset from the phone in the lab and dialed out. A gruff voice
  acknowledged me. I launched into my rehearsed pitch. ``I've just graduated
  with a PhD in theoretical Computer Science. '' He cut me off. ``We want
  workers, not thinkers.'' and the line went dead. 

  ``Workers not thinkers'' --
  the phrase has haunted me these many years. Maybe the shoe is finally on the
  other foot?
  \ee

\subsection{Future Work}

We are sharing this work before all the rough edges have been smoothed 
out in the hope of getting feedback. We would especially like suggestions on
\be
\item packaging, versioning, \ldots
\item extending code generation to other parts of the system, such as unit
  testing
\item cross-testing different LLMs against each other
\item viable replacements for the role C and Lua played in this implementation
  \ee

\bibliographystyle{alpha}
\bibliography{ref}
\appendix

\input{doc_qdf_definition}

\end{document}

